% =============================================================================
% BMA Health — Preprocessing Pipeline Technical Note (XeLaTeX via Tectonic)
% =============================================================================
% Compile from api/templates/latex/ so that `Path=assets/` resolves:
%     cd api/templates/latex && tectonic preprocessing.tex
% =============================================================================
\documentclass[11pt,a4paper]{article}

\newcommand{\bmaheaderleft}{BMA Health --- Preprocessing Pipeline}
\newcommand{\bmaversion}{v1.0 · 2026-04-22}
% ============================================================
% BMA Health Article Preamble (XeLaTeX via Tectonic)
% ============================================================
% Before % ============================================================
% BMA Health Article Preamble (XeLaTeX via Tectonic)
% ============================================================
% Before % ============================================================
% BMA Health Article Preamble (XeLaTeX via Tectonic)
% ============================================================
% Before \input{bma_article_preamble}, define:
%   \bmaheaderleft   e.g. {BMA Health --- Screening Report}
%   \bmaversion      e.g. {Report v1.0}
% ============================================================

% --- Packages ---
\usepackage{geometry}
\geometry{a4paper, margin=2.0cm}
\usepackage{hyperref}
\usepackage{parskip}
\usepackage{booktabs}
\usepackage{amsmath,amssymb}
\usepackage{graphicx}
\usepackage{xcolor}
\usepackage{float}
\usepackage{needspace}
\usepackage{caption}
\usepackage{fancyhdr}
\usepackage{enumitem}
\usepackage{colortbl}
\usepackage{multirow}
\usepackage{array}
\usepackage{pifont}
\usepackage{longtable}
\usepackage{fontspec}
\usepackage{polyglossia}

% --- Thai line breaking (critical for word wrap) ---
\XeTeXlinebreaklocale "th"
\XeTeXlinebreakskip = 0pt plus 1pt
\emergencystretch=3em
\tolerance=9999
\usepackage{lastpage}

% S7 fix: ChartBlock.render_latex emits raw \begin{tikzpicture}/pgfplots
% but no upstream preamble loaded the packages. Adding here so descriptor-
% driven reports compile when they include any chart block.
\usepackage{tikz}
\usepackage{pgfplots}
\pgfplotsset{compat=1.18}

% --- Colors (BMA Green palette) ---
\definecolor{bmagreen}{HTML}{00744B}
\definecolor{bmagreenlight}{HTML}{4CAF50}
\definecolor{bmadark}{HTML}{005F3D}
\definecolor{errred}{HTML}{D32F2F}
\definecolor{okgreen}{HTML}{388E3C}
\definecolor{warnblue}{HTML}{1565C0}
\definecolor{warnamber}{HTML}{F57F17}
\definecolor{lightgray}{gray}{0.95}
\definecolor{tablegray}{gray}{0.97}

% --- Header/Footer ---
% S10: ``\pageref{LastPage}`` requires ``\usepackage{lastpage}`` (loaded
% above on line 36). Footer reads ``1 / 12`` so readers can see total pages.
\pagestyle{fancy}
\fancyhf{}
\fancyhead[L]{\small\color{bmagreen}\bmafont\textbf{\bmaheaderleft}}
\fancyhead[R]{\small\color{gray}\bmaversion}
\fancyfoot[C]{\small\thepage\ / \pageref{LastPage}}
\renewcommand{\headrulewidth}{0.4pt}
\renewcommand{\footrulewidth}{0.0pt}
\renewcommand{\headrule}{\hbox to\headwidth{\color{bmagreen}\leaders\hrule height \headrulewidth\hfill}}

% Plain pages (titlepage / chapter starts) get nothing — keeps the cover
% clean of fancy header/footer chrome.
\fancypagestyle{plain}{%
  \fancyhf{}%
  \renewcommand{\headrulewidth}{0pt}%
}

% --- Section color ---
\usepackage{titlesec}
\titleformat{\section}{\large\bmafont\bfseries\color{bmagreen}}{\thesection}{1em}{}
\titleformat{\subsection}{\normalsize\bmafont\bfseries\color{bmagreen!80!black}}{\thesubsection}{1em}{}
\titleformat{\subsubsection}{\normalsize\bfseries}{\thesubsubsection}{1em}{}

% --- Table of Contents styling ---
\usepackage{tocloft}
\renewcommand{\cftsecfont}{\bfseries\color{bmagreen}}
\renewcommand{\cftsecpagefont}{\bfseries\color{bmagreen}}
\renewcommand{\cftsubsecfont}{\color{bmadark}}
\renewcommand{\cftsubsecpagefont}{\color{bmadark}}

% --- Caption style ---
\captionsetup{font=small, labelfont={bf,color=bmagreen}}

% --- Caption-convention helpers (Phase 0) ---
% Always use these so figures get caption BELOW, tables get caption ABOVE
% (academic / scientific convention). Mirrors helpers in
% ``services/reports/blocks/_render_helpers.py``.
%
%   \bmafig[width]{path}{caption}{label}  -- figure, caption below
%   \bmatab{caption}{label}{tabular_body} -- table, caption above
%
% Default width 0.85\textwidth keeps figures readable without bleeding
% into the margin. Override with \bmafig[\textwidth]{...}{...}{...} for
% full-page figures.
\newcommand{\bmafig}[4][0.85\textwidth]{%
  \begin{figure}[H]%
    \centering
    \includegraphics[width=#1]{#2}%
    \caption{#3}\label{#4}%
  \end{figure}%
}
\newcommand{\bmatab}[3]{%
  \begin{table}[H]%
    \centering
    \caption{#1}\label{#2}%
    #3%
  \end{table}%
}

% --- Hyperref: colors + PDF metadata + bookmark tree (S10) ---
% PDF readers (Preview, Acrobat) show the metadata in their info panel
% and the bookmark tree as a navigable side panel matching the TOC.
\hypersetup{
  pdftitle={\bmaheaderleft},
  pdfauthor={BMA Medical Service Department},
  pdfsubject={Bangkok population health screening report},
  pdfkeywords={NCDs, screening, Bangkok, BMA, health zone},
  colorlinks=true,
  linkcolor=bmagreen,
  urlcolor=bmagreen!70!black,
  citecolor=bmagreen,
  bookmarksopen=true,
  bookmarksnumbered=true
}

% --- Font: Google Sans (official, Thai+Latin) ---
\newfontfamily\bodyfont{GoogleSans-Regular.ttf}[
  Path=assets/,
  BoldFont=GoogleSans-Bold.ttf,
]
% \bmafont alias for backward compat with templates
\let\bmafont\bodyfont

% --- Convenience commands ---
% Risk cell coloring: \riskcell{45.3}
\newcommand{\riskcell}[1]{%
  \ifdim #1pt > 40pt
    \cellcolor{errred!15}\textcolor{errred}{\textbf{#1\%}}%
  \else\ifdim #1pt > 25pt
    \cellcolor{warnblue!10}\textcolor{warnblue}{#1\%}%
  \else
    \cellcolor{okgreen!10}\textcolor{okgreen}{#1\%}%
  \fi\fi
}

% Trend arrows
\newcommand{\trendUp}{\textcolor{errred}{$\blacktriangle$}}
\newcommand{\trendDown}{\textcolor{okgreen}{$\blacktriangledown$}}
\newcommand{\trendStable}{\textcolor{gray}{$\bullet$}}

% Statistical significance marker
\newcommand{\sigmark}[1]{%
  \ifdim #1pt < 0.001pt
    \textcolor{errred}{\textbf{***}}%
  \else\ifdim #1pt < 0.01pt
    \textcolor{errred}{\textbf{**}}%
  \else\ifdim #1pt < 0.05pt
    \textcolor{warnblue}{\textbf{*}}%
  \else
    \textcolor{gray}{ns}%
  \fi\fi\fi
}

% Compact list
\newlist{compactlist}{itemize}{3}
\setlist[compactlist]{
  topsep=2pt,
  partopsep=0pt,
  itemsep=1pt,
  parsep=0pt,
  leftmargin=1.5em,
  label=\textcolor{bmagreen}{\textbullet},
}

% Key-value row for summary boxes
\newcommand{\kvrow}[2]{%
  \textcolor{bmadark}{\textbf{#1}} & #2 \\
}
, define:
%   \bmaheaderleft   e.g. {BMA Health --- Screening Report}
%   \bmaversion      e.g. {Report v1.0}
% ============================================================

% --- Packages ---
\usepackage{geometry}
\geometry{a4paper, margin=2.0cm}
\usepackage{hyperref}
\usepackage{parskip}
\usepackage{booktabs}
\usepackage{amsmath,amssymb}
\usepackage{graphicx}
\usepackage{xcolor}
\usepackage{float}
\usepackage{needspace}
\usepackage{caption}
\usepackage{fancyhdr}
\usepackage{enumitem}
\usepackage{colortbl}
\usepackage{multirow}
\usepackage{array}
\usepackage{pifont}
\usepackage{longtable}
\usepackage{fontspec}
\usepackage{polyglossia}

% --- Thai line breaking (critical for word wrap) ---
\XeTeXlinebreaklocale "th"
\XeTeXlinebreakskip = 0pt plus 1pt
\emergencystretch=3em
\tolerance=9999
\usepackage{lastpage}

% S7 fix: ChartBlock.render_latex emits raw \begin{tikzpicture}/pgfplots
% but no upstream preamble loaded the packages. Adding here so descriptor-
% driven reports compile when they include any chart block.
\usepackage{tikz}
\usepackage{pgfplots}
\pgfplotsset{compat=1.18}

% --- Colors (BMA Green palette) ---
\definecolor{bmagreen}{HTML}{00744B}
\definecolor{bmagreenlight}{HTML}{4CAF50}
\definecolor{bmadark}{HTML}{005F3D}
\definecolor{errred}{HTML}{D32F2F}
\definecolor{okgreen}{HTML}{388E3C}
\definecolor{warnblue}{HTML}{1565C0}
\definecolor{warnamber}{HTML}{F57F17}
\definecolor{lightgray}{gray}{0.95}
\definecolor{tablegray}{gray}{0.97}

% --- Header/Footer ---
% S10: ``\pageref{LastPage}`` requires ``\usepackage{lastpage}`` (loaded
% above on line 36). Footer reads ``1 / 12`` so readers can see total pages.
\pagestyle{fancy}
\fancyhf{}
\fancyhead[L]{\small\color{bmagreen}\bmafont\textbf{\bmaheaderleft}}
\fancyhead[R]{\small\color{gray}\bmaversion}
\fancyfoot[C]{\small\thepage\ / \pageref{LastPage}}
\renewcommand{\headrulewidth}{0.4pt}
\renewcommand{\footrulewidth}{0.0pt}
\renewcommand{\headrule}{\hbox to\headwidth{\color{bmagreen}\leaders\hrule height \headrulewidth\hfill}}

% Plain pages (titlepage / chapter starts) get nothing — keeps the cover
% clean of fancy header/footer chrome.
\fancypagestyle{plain}{%
  \fancyhf{}%
  \renewcommand{\headrulewidth}{0pt}%
}

% --- Section color ---
\usepackage{titlesec}
\titleformat{\section}{\large\bmafont\bfseries\color{bmagreen}}{\thesection}{1em}{}
\titleformat{\subsection}{\normalsize\bmafont\bfseries\color{bmagreen!80!black}}{\thesubsection}{1em}{}
\titleformat{\subsubsection}{\normalsize\bfseries}{\thesubsubsection}{1em}{}

% --- Table of Contents styling ---
\usepackage{tocloft}
\renewcommand{\cftsecfont}{\bfseries\color{bmagreen}}
\renewcommand{\cftsecpagefont}{\bfseries\color{bmagreen}}
\renewcommand{\cftsubsecfont}{\color{bmadark}}
\renewcommand{\cftsubsecpagefont}{\color{bmadark}}

% --- Caption style ---
\captionsetup{font=small, labelfont={bf,color=bmagreen}}

% --- Caption-convention helpers (Phase 0) ---
% Always use these so figures get caption BELOW, tables get caption ABOVE
% (academic / scientific convention). Mirrors helpers in
% ``services/reports/blocks/_render_helpers.py``.
%
%   \bmafig[width]{path}{caption}{label}  -- figure, caption below
%   \bmatab{caption}{label}{tabular_body} -- table, caption above
%
% Default width 0.85\textwidth keeps figures readable without bleeding
% into the margin. Override with \bmafig[\textwidth]{...}{...}{...} for
% full-page figures.
\newcommand{\bmafig}[4][0.85\textwidth]{%
  \begin{figure}[H]%
    \centering
    \includegraphics[width=#1]{#2}%
    \caption{#3}\label{#4}%
  \end{figure}%
}
\newcommand{\bmatab}[3]{%
  \begin{table}[H]%
    \centering
    \caption{#1}\label{#2}%
    #3%
  \end{table}%
}

% --- Hyperref: colors + PDF metadata + bookmark tree (S10) ---
% PDF readers (Preview, Acrobat) show the metadata in their info panel
% and the bookmark tree as a navigable side panel matching the TOC.
\hypersetup{
  pdftitle={\bmaheaderleft},
  pdfauthor={BMA Medical Service Department},
  pdfsubject={Bangkok population health screening report},
  pdfkeywords={NCDs, screening, Bangkok, BMA, health zone},
  colorlinks=true,
  linkcolor=bmagreen,
  urlcolor=bmagreen!70!black,
  citecolor=bmagreen,
  bookmarksopen=true,
  bookmarksnumbered=true
}

% --- Font: Google Sans (official, Thai+Latin) ---
\newfontfamily\bodyfont{GoogleSans-Regular.ttf}[
  Path=assets/,
  BoldFont=GoogleSans-Bold.ttf,
]
% \bmafont alias for backward compat with templates
\let\bmafont\bodyfont

% --- Convenience commands ---
% Risk cell coloring: \riskcell{45.3}
\newcommand{\riskcell}[1]{%
  \ifdim #1pt > 40pt
    \cellcolor{errred!15}\textcolor{errred}{\textbf{#1\%}}%
  \else\ifdim #1pt > 25pt
    \cellcolor{warnblue!10}\textcolor{warnblue}{#1\%}%
  \else
    \cellcolor{okgreen!10}\textcolor{okgreen}{#1\%}%
  \fi\fi
}

% Trend arrows
\newcommand{\trendUp}{\textcolor{errred}{$\blacktriangle$}}
\newcommand{\trendDown}{\textcolor{okgreen}{$\blacktriangledown$}}
\newcommand{\trendStable}{\textcolor{gray}{$\bullet$}}

% Statistical significance marker
\newcommand{\sigmark}[1]{%
  \ifdim #1pt < 0.001pt
    \textcolor{errred}{\textbf{***}}%
  \else\ifdim #1pt < 0.01pt
    \textcolor{errred}{\textbf{**}}%
  \else\ifdim #1pt < 0.05pt
    \textcolor{warnblue}{\textbf{*}}%
  \else
    \textcolor{gray}{ns}%
  \fi\fi\fi
}

% Compact list
\newlist{compactlist}{itemize}{3}
\setlist[compactlist]{
  topsep=2pt,
  partopsep=0pt,
  itemsep=1pt,
  parsep=0pt,
  leftmargin=1.5em,
  label=\textcolor{bmagreen}{\textbullet},
}

% Key-value row for summary boxes
\newcommand{\kvrow}[2]{%
  \textcolor{bmadark}{\textbf{#1}} & #2 \\
}
, define:
%   \bmaheaderleft   e.g. {BMA Health --- Screening Report}
%   \bmaversion      e.g. {Report v1.0}
% ============================================================

% --- Packages ---
\usepackage{geometry}
\geometry{a4paper, margin=2.0cm}
\usepackage{hyperref}
\usepackage{parskip}
\usepackage{booktabs}
\usepackage{amsmath,amssymb}
\usepackage{graphicx}
\usepackage{xcolor}
\usepackage{float}
\usepackage{needspace}
\usepackage{caption}
\usepackage{fancyhdr}
\usepackage{enumitem}
\usepackage{colortbl}
\usepackage{multirow}
\usepackage{array}
\usepackage{pifont}
\usepackage{longtable}
\usepackage{fontspec}
\usepackage{polyglossia}

% --- Thai line breaking (critical for word wrap) ---
\XeTeXlinebreaklocale "th"
\XeTeXlinebreakskip = 0pt plus 1pt
\emergencystretch=3em
\tolerance=9999
\usepackage{lastpage}

% S7 fix: ChartBlock.render_latex emits raw \begin{tikzpicture}/pgfplots
% but no upstream preamble loaded the packages. Adding here so descriptor-
% driven reports compile when they include any chart block.
\usepackage{tikz}
\usepackage{pgfplots}
\pgfplotsset{compat=1.18}

% --- Colors (BMA Green palette) ---
\definecolor{bmagreen}{HTML}{00744B}
\definecolor{bmagreenlight}{HTML}{4CAF50}
\definecolor{bmadark}{HTML}{005F3D}
\definecolor{errred}{HTML}{D32F2F}
\definecolor{okgreen}{HTML}{388E3C}
\definecolor{warnblue}{HTML}{1565C0}
\definecolor{warnamber}{HTML}{F57F17}
\definecolor{lightgray}{gray}{0.95}
\definecolor{tablegray}{gray}{0.97}

% --- Header/Footer ---
% S10: ``\pageref{LastPage}`` requires ``\usepackage{lastpage}`` (loaded
% above on line 36). Footer reads ``1 / 12`` so readers can see total pages.
\pagestyle{fancy}
\fancyhf{}
\fancyhead[L]{\small\color{bmagreen}\bmafont\textbf{\bmaheaderleft}}
\fancyhead[R]{\small\color{gray}\bmaversion}
\fancyfoot[C]{\small\thepage\ / \pageref{LastPage}}
\renewcommand{\headrulewidth}{0.4pt}
\renewcommand{\footrulewidth}{0.0pt}
\renewcommand{\headrule}{\hbox to\headwidth{\color{bmagreen}\leaders\hrule height \headrulewidth\hfill}}

% Plain pages (titlepage / chapter starts) get nothing — keeps the cover
% clean of fancy header/footer chrome.
\fancypagestyle{plain}{%
  \fancyhf{}%
  \renewcommand{\headrulewidth}{0pt}%
}

% --- Section color ---
\usepackage{titlesec}
\titleformat{\section}{\large\bmafont\bfseries\color{bmagreen}}{\thesection}{1em}{}
\titleformat{\subsection}{\normalsize\bmafont\bfseries\color{bmagreen!80!black}}{\thesubsection}{1em}{}
\titleformat{\subsubsection}{\normalsize\bfseries}{\thesubsubsection}{1em}{}

% --- Table of Contents styling ---
\usepackage{tocloft}
\renewcommand{\cftsecfont}{\bfseries\color{bmagreen}}
\renewcommand{\cftsecpagefont}{\bfseries\color{bmagreen}}
\renewcommand{\cftsubsecfont}{\color{bmadark}}
\renewcommand{\cftsubsecpagefont}{\color{bmadark}}

% --- Caption style ---
\captionsetup{font=small, labelfont={bf,color=bmagreen}}

% --- Caption-convention helpers (Phase 0) ---
% Always use these so figures get caption BELOW, tables get caption ABOVE
% (academic / scientific convention). Mirrors helpers in
% ``services/reports/blocks/_render_helpers.py``.
%
%   \bmafig[width]{path}{caption}{label}  -- figure, caption below
%   \bmatab{caption}{label}{tabular_body} -- table, caption above
%
% Default width 0.85\textwidth keeps figures readable without bleeding
% into the margin. Override with \bmafig[\textwidth]{...}{...}{...} for
% full-page figures.
\newcommand{\bmafig}[4][0.85\textwidth]{%
  \begin{figure}[H]%
    \centering
    \includegraphics[width=#1]{#2}%
    \caption{#3}\label{#4}%
  \end{figure}%
}
\newcommand{\bmatab}[3]{%
  \begin{table}[H]%
    \centering
    \caption{#1}\label{#2}%
    #3%
  \end{table}%
}

% --- Hyperref: colors + PDF metadata + bookmark tree (S10) ---
% PDF readers (Preview, Acrobat) show the metadata in their info panel
% and the bookmark tree as a navigable side panel matching the TOC.
\hypersetup{
  pdftitle={\bmaheaderleft},
  pdfauthor={BMA Medical Service Department},
  pdfsubject={Bangkok population health screening report},
  pdfkeywords={NCDs, screening, Bangkok, BMA, health zone},
  colorlinks=true,
  linkcolor=bmagreen,
  urlcolor=bmagreen!70!black,
  citecolor=bmagreen,
  bookmarksopen=true,
  bookmarksnumbered=true
}

% --- Font: Google Sans (official, Thai+Latin) ---
\newfontfamily\bodyfont{GoogleSans-Regular.ttf}[
  Path=assets/,
  BoldFont=GoogleSans-Bold.ttf,
]
% \bmafont alias for backward compat with templates
\let\bmafont\bodyfont

% --- Convenience commands ---
% Risk cell coloring: \riskcell{45.3}
\newcommand{\riskcell}[1]{%
  \ifdim #1pt > 40pt
    \cellcolor{errred!15}\textcolor{errred}{\textbf{#1\%}}%
  \else\ifdim #1pt > 25pt
    \cellcolor{warnblue!10}\textcolor{warnblue}{#1\%}%
  \else
    \cellcolor{okgreen!10}\textcolor{okgreen}{#1\%}%
  \fi\fi
}

% Trend arrows
\newcommand{\trendUp}{\textcolor{errred}{$\blacktriangle$}}
\newcommand{\trendDown}{\textcolor{okgreen}{$\blacktriangledown$}}
\newcommand{\trendStable}{\textcolor{gray}{$\bullet$}}

% Statistical significance marker
\newcommand{\sigmark}[1]{%
  \ifdim #1pt < 0.001pt
    \textcolor{errred}{\textbf{***}}%
  \else\ifdim #1pt < 0.01pt
    \textcolor{errred}{\textbf{**}}%
  \else\ifdim #1pt < 0.05pt
    \textcolor{warnblue}{\textbf{*}}%
  \else
    \textcolor{gray}{ns}%
  \fi\fi\fi
}

% Compact list
\newlist{compactlist}{itemize}{3}
\setlist[compactlist]{
  topsep=2pt,
  partopsep=0pt,
  itemsep=1pt,
  parsep=0pt,
  leftmargin=1.5em,
  label=\textcolor{bmagreen}{\textbullet},
}

% Key-value row for summary boxes
\newcommand{\kvrow}[2]{%
  \textcolor{bmadark}{\textbf{#1}} & #2 \\
}


% --- Thai font (Google Sans has native Thai support) ---
\setmainfont{GoogleSans-Regular.ttf}[
  Path=assets/,
  BoldFont=GoogleSans-Bold.ttf,
]
\setsansfont{GoogleSans-Regular.ttf}[
  Path=assets/,
  BoldFont=GoogleSans-Bold.ttf,
]
% --- Monospace: use a system font that supports Thai so \texttt{} with Thai
%     characters renders correctly (fallback from LaTeX's default lmmono).
\setmonofont{Menlo}[Scale=0.88]

% --- listings for code samples ---
\usepackage{listings}
\definecolor{codebg}{HTML}{F6F8F7}
\definecolor{codekey}{HTML}{00744B}
\definecolor{codestr}{HTML}{B7410E}
\definecolor{codecomment}{HTML}{718096}

% Custom "monospace-ish" font that supports Thai (uses the main Google Sans,
% which has Thai glyphs, just at slightly smaller size + no proportional
% kerning). listings' basicstyle can take any font family.
\newfontfamily\codefont{GoogleSans-Regular.ttf}[
  Path=assets/,
  Scale=0.88,
  BoldFont=GoogleSans-Bold.ttf,
]

\lstdefinestyle{shellstyle}{
  basicstyle=\codefont\small,
  backgroundcolor=\color{codebg},
  frame=single,
  framerule=0.3pt,
  rulecolor=\color{gray!40},
  breaklines=true,
  breakatwhitespace=true,
  columns=fullflexible,
  keepspaces=true,
  xleftmargin=8pt,
  xrightmargin=4pt,
  lineskip=1pt,
  aboveskip=6pt,
  belowskip=6pt,
  captionpos=b,
}
\lstdefinestyle{sqlstyle}{
  style=shellstyle,
  language=SQL,
  keywordstyle=\color{codekey}\bfseries,
  stringstyle=\color{codestr},
  commentstyle=\color{codecomment}\itshape,
  morekeywords={CASCADE,CONCURRENTLY,RETURNING,COALESCE,FILTER,STRING_AGG}
}
\lstdefinestyle{pythonstyle}{
  style=shellstyle,
  language=Python,
  keywordstyle=\color{codekey}\bfseries,
  stringstyle=\color{codestr},
  commentstyle=\color{codecomment}\itshape,
}

% Let tables break across pages automatically
\usepackage{tabularx}
\newcolumntype{L}[1]{>{\raggedright\arraybackslash}p{#1}}

% Since the preamble loads `fancyhdr` and expects \bmaheaderleft, the header
% will already appear on every page. No \maketitle, no cover.

\begin{document}

% No title page — go straight to content. A compact heading bar instead.
\begin{center}
  \vspace*{-0.6cm}
  \textcolor{bmagreen}{\rule{\linewidth}{1.2pt}}\\[2pt]
  {\Large\bfseries\color{bmagreen} BMA Health Preprocessing Pipeline}\\[2pt]
  {\small\color{gray} โครงการคัดกรองสุขภาพกรุงเทพมหานคร · สำนักการแพทย์ กรุงเทพมหานคร}\\[2pt]
  \textcolor{bmagreen}{\rule{\linewidth}{1.2pt}}
\end{center}

\vspace{0.3cm}

\noindent
เอกสารนี้สรุป \textbf{preprocessing pipeline ทั้งหมด} ที่แปลงไฟล์ CSV ต้นทาง 13 ไฟล์จาก 3 ระบบ
(Portal / App1 / App2) ให้กลายเป็นตาราง PostgreSQL และ materialized views สำหรับ API และรายงาน.
ไม่มีการ impute ค่าใด ๆ — ค่าผิด/เกินช่วง/ว่าง จะถูกแปลงเป็น \texttt{NULL}.
ข้อมูลคนเดียวกันที่ปรากฏในหลายระบบจะ \textit{stack} เก็บเป็นหลายแถว (ไม่ merge) เพื่อรักษาจำนวนการตรวจจริง.

\tableofcontents
\vspace{0.5cm}

% =============================================================================
\section{หลักการออกแบบ 5 ข้อ}
% =============================================================================

\begin{enumerate}[leftmargin=*, itemsep=3pt]
  \item \textbf{No imputation} --- ค่าผิดพลาด/นอกช่วง/ว่างเปล่า $\rightarrow$ \texttt{NULL} เสมอ.
    ไม่เดา ไม่เติมด้วยค่าเฉลี่ย ไม่ใช้ last-observation-carried-forward.
    Imputation บิดเบือนการวิเคราะห์ epidemiology.
  \item \textbf{Stack mode} --- คนเดียวกันปรากฏใน 2 ระบบ $=$ สร้าง 2 แถวใน \texttt{raw\_patients}
    (แถวละ \texttt{data\_source} ต่างกัน) เพื่อรักษาจำนวนการตรวจจริงและรายงานความซ้ำซ้อนข้ามระบบ.
  \item \textbf{3-tier fields} --- \textit{raw} / \textit{source-derived} (\texttt{*\_src}) / \textit{ETL-derived}
    แยกกันชัดเจน ตรวจสอบย้อนกลับได้และเปรียบเทียบระหว่าง source ได้.
  \item \textbf{Provenance} --- ทุก record มี \texttt{data\_source} (\texttt{portal}/\texttt{app1}/\texttt{app2})
    และ \texttt{import\_batch\_id} เก็บไว้เสมอ.
  \item \textbf{Fail-soft} --- \texttt{execute\_values} มี \texttt{SAVEPOINT} แถว-ต่อ-แถว.
    แถวที่พังจะถูก skip พร้อม log. ไม่ทำให้ transaction ทั้งก้อนล้ม.
\end{enumerate}

% =============================================================================
\section{แหล่งข้อมูล 3 ระบบ}
% =============================================================================

\subsection{โครงสร้างโฟลเดอร์}

Web upload และ CLI คาดหวังโครงสร้างดังนี้:
\begin{lstlisting}[style=shellstyle]
<data-dir>/
|- portal/   (7 files) pt.csv, pthistory.csv, vitalsignslf.csv,
|                      homevisit.csv, homehealth.csv,
|                      labhealth.csv, labhealthext.csv
|- app1/     (5 files) pt.csv, vitalsignslf.csv, homevisit.csv,
|                      homehealth.csv, labhealth.csv
|- app2/     (1 file)  app2.csv
\end{lstlisting}

\subsection{ความแตกต่างระหว่าง source}

\begin{center}
\small
\begin{tabularx}{\linewidth}{l X X X}
\toprule
& \textbf{Portal} & \textbf{App1} & \textbf{App2} \\
\midrule
จำนวนไฟล์          & 7 & 5 & 1 \\
Columns รวม        & 383 & 189 & 103 \\
ID field           & \texttt{IDCARD} & \texttt{PID} & \texttt{PID} \\
Birth date         & \texttt{BIRTHDATE} & \texttt{BRTHDATE} & ไม่มี (มีแค่ \texttt{AGE\_GROUP}) \\
Sex encoding       & 10 / 20 & 10 / 20 & "ชาย" / "หญิง" \\
BMI                & ต้องคำนวณเอง & ต้องคำนวณเอง & pre-computed ใน \texttt{BMI} \\
Lab results        & ตัวเลข raw (HMGB, FBS, CHOLEST) & เหมือน Portal & ตัวเลข 3 ช่อง + text interpret 11 ช่อง \\
Screening flags    & 0/1 integer & 0/1 integer & "ปกติ" / "เสี่ยง" / "ผิดปกติ" \\
\bottomrule
\end{tabularx}
\end{center}

% =============================================================================
\section{Pipeline 6 ขั้นตอน}
% =============================================================================

\begin{center}
\footnotesize
\begin{tabular}{rlL{8.5cm}}
\toprule
\# & ขั้นตอน & คำอธิบาย \\
\midrule
1 & Upload (streaming)  & \texttt{shutil.copyfileobj} ลง tempfile บน disk — 64 KB chunks.
  Peak memory $< 10$ MB แม้ไฟล์ 500 MB. \\
2 & Parse               & \texttt{pd.read\_csv(path, dtype=str, low\_memory=False)}.
  Encoding fallback: UTF-8 $\rightarrow$ TIS-620 $\rightarrow$ CP874 $\rightarrow$ Latin-1. \\
3 & Hash                & Base64 decode $\rightarrow$ HMAC-SHA-256 ด้วย env secret $\rightarrow$ hex digest 64 chars.
  Portal/App1/App2 ใช้ secret เดียวกัน $\rightarrow$ match ข้าม source ได้. \\
4 & Clean               & Out-of-range $\rightarrow$ \texttt{NULL} (ดู \S\ref{sec:cleanse}). \\
5 & Stack insert        & Plain \texttt{INSERT} (ไม่มี \texttt{ON CONFLICT} บน child tables).
  Natural key = \texttt{(idcard\_hash, data\_source)}. \\
6 & View refresh        & \texttt{REFRESH MATERIALIZED VIEW CONCURRENTLY} ทั้ง 13 views.
  Non-fatal — ถ้า refresh ล้ม ข้อมูล raw ยัง committed แล้ว. \\
\bottomrule
\end{tabular}
\end{center}

% =============================================================================
\section{Data Cleansing Rules}
\label{sec:cleanse}
% =============================================================================

\subsection{Filter out (ลบทิ้ง)}

\begin{itemize}[leftmargin=*, itemsep=1pt]
  \item \texttt{IDCARD}/\texttt{PID} เป็น \texttt{NULL}, ว่าง, หรือ base64 decode ไม่ได้
    $\rightarrow$ ลบ (ระบุตัวบุคคลไม่ได้, $< 1$\%).
\end{itemize}

\noindent เราไม่ลบ row ที่ \texttt{CANCELST}$=1$ อีกต่อไป — เก็บไว้
และใช้ \texttt{cancel\_status} filter ตอน query เพราะยังต้องรายงาน.

\subsection{Out-of-range $\rightarrow$ NULL}

หลักการ: ค่าที่ \textit{เป็นไปไม่ได้ทาง clinical} $\rightarrow$ \texttt{NULL} --- ไม่ใช่ 0 ไม่ใช่ค่าเฉลี่ย.

\begin{center}
\small
\begin{tabularx}{\linewidth}{l l X}
\toprule
\textbf{Field} & \textbf{เงื่อนไข NULL} & \textbf{เหตุผล} \\
\midrule
\texttt{BIRTHDATE} ปี & $< 1900$ หรือ $> 2030$ & อายุสูงสุดโลก $= 122$ ปี \\
\texttt{age}          & $< 0$ หรือ $> 150$ & เป็นไปไม่ได้ \\
\texttt{HEIGHT}       & $< 50$ หรือ $> 250$ cm & ทารก $\sim 50$, คนสูงสุด $272$ cm \\
\texttt{WEIGHT}       & $< 10$ หรือ $> 300$ kg & ทารก $\sim 3$, คนหนักสุดในไทย $\sim 300$ kg \\
\texttt{WSTL}         & $< 30$ cm & \textbf{พบจริง} มีค่า $0$ (ไม่ได้วัดแต่กรอก 0) \\
\texttt{HBPN} (SBP)   & $< 40$ หรือ $> 300$ mmHg & \textbf{พบจริง} มีค่า $0, 1$ (ไม่ได้วัด) \\
\texttt{LBPN} (DBP)   & $< 20$ หรือ $> 200$ mmHg & DBP $< 20 =$ shock รุนแรง \\
\texttt{FBS, CHOLEST} & $< 0$ หรือ $> 999$ mg/dL & เครื่องวัดสูงสุด \\
\texttt{HMGB}         & $< 0$ หรือ $> 30$ g/dL & Hb สูงสุดจริง $\sim 20$ (polycythemia vera) \\
\texttt{EGFR}         & $< 0$ หรือ $> 200$ mL/min & Clinical range \\
\texttt{BMI} (computed) & $> 80$ & คนหนักสุดโลก BMI $\sim 70$ \\
integer               & $> 2{,}147{,}483{,}647$ & เกิน PostgreSQL INT4 \\
float                 & $\infty$ หรือ NaN & data เสีย \\
\bottomrule
\end{tabularx}
\end{center}

\subsection{Transform (แปลงค่า)}

\begin{center}
\footnotesize
\begin{tabularx}{\linewidth}{L{6cm} X}
\toprule
\textbf{เงื่อนไข} & \textbf{การแปลง} \\
\midrule
\texttt{BIRTHDATE} ปี $> 2400$     & ลบ 543 (พ.ศ.$\rightarrow$ค.ศ.) \\
App2 \texttt{MALE} = "ชาย"/"หญิง"   & $\rightarrow$ 10 / 20 \\
App2 \texttt{DM\_NAME} = "ปกติ"/"เสี่ยง"/"เป็น" & $\rightarrow$ 0 / 1 / 1 \\
App2 \texttt{H\_DM\_NAME} = "เป็น"/"ไม่เป็น"/"ไม่แน่ใจ" & $\rightarrow$ True / False / NULL \\
App2 \texttt{DISTRICT} $= 9999$    & $\rightarrow$ \texttt{NULL} (placeholder "ไม่ระบุ") \\
\texttt{DISTRICTBKK} ว่าง          & Backfill 4-priority (ดู \S\ref{sec:backfill}) \\
\bottomrule
\end{tabularx}
\end{center}

\subsection{District backfill (4 priority)}
\label{sec:backfill}

\begin{enumerate}[leftmargin=*, itemsep=1pt]
  \item \texttt{DISTRICTBKK} จาก \texttt{vitalsignslf.csv} (ถ้ามี)
  \item \texttt{DISTRICT} จาก \texttt{app2.csv} (ถ้ามี, $\ne 9999$)
  \item \texttt{facility\_code} $\rightarrow$ \texttt{ref\_facility\_districts} mapping
  \item \texttt{home\_district} จาก \texttt{homevisit} (ต้องอยู่ $1001$--$1050$)
\end{enumerate}

\subsection{Row-level error handling}

\begin{lstlisting}[style=pythonstyle]
execute_values(cur, sql, rows, page_size=2000, progress_cb=cb):
  for chunk in chunks(rows, page_size):
    SAVEPOINT ev_batch
    try:
      INSERT chunk   # batch-mode (fast path)
      RELEASE SAVEPOINT
    except:
      ROLLBACK TO SAVEPOINT ev_batch
      # retry each row; skip offenders
      for row in chunk:
        try:    SAVEPOINT ev_row; INSERT row; RELEASE
        except: ROLLBACK TO SAVEPOINT; skipped += 1
    progress_cb(done, total)   # UI live progress
\end{lstlisting}

% =============================================================================
\section{Stack Mode --- ไม่ Merge ไม่ Dedupe}
% =============================================================================

\subsection{Natural key ใหม่}

\begin{lstlisting}[style=sqlstyle]
-- Before (migration 001-004):
raw_patients (idcard_hash UNIQUE, ...)

-- After (migration 011):
raw_patients (idcard_hash, data_source, ...)
  UNIQUE(idcard_hash, data_source)
-- Same person in 2 systems = 2 rows (one per source).
\end{lstlisting}

ตาราง child (\texttt{raw\_visits}, \texttt{raw\_vitalsigns}, \ldots) มี FK \texttt{patient\_id}
ชี้ไปที่ \texttt{raw\_patients(id)} ของ source ตัวเอง.
Composite \texttt{UNIQUE(patient\_id, visit\_date, facility\_code)} ถูกถอดเหลือเป็น \texttt{INDEX}
(visit เดียวกันจาก 2 ระบบ = ให้เป็น 2 row ได้, นับเป็นการตรวจซ้ำ).

\subsection{ผลต่อ query}

\begin{lstlisting}[style=sqlstyle]
-- นับจำนวนการตรวจ (visits, รวมซ้ำข้ามระบบ):
SELECT COUNT(*) FROM raw_vitalsigns;

-- นับจำนวนประชาชนเฉพาะ (unique people):
SELECT COUNT(DISTINCT idcard_hash) FROM raw_patients;

-- ตรวจซ้ำข้ามระบบ (view built-in):
SELECT * FROM v_cross_system_duplicates;
-- columns: idcard_hash, n_systems, systems, n_patient_rows

-- สรุปจำนวน rows แต่ละ source แต่ละ table:
SELECT * FROM v_source_row_counts;
\end{lstlisting}

\subsection{Best-available value}

\begin{lstlisting}[style=sqlstyle]
-- BMI ที่ดีที่สุด: priority Portal calc > App1 calc > App2 src
SELECT idcard_hash,
  COALESCE(
    MAX(bmi)     FILTER (WHERE data_source = 'portal'),
    MAX(bmi)     FILTER (WHERE data_source = 'app1'),
    MAX(bmi_src) FILTER (WHERE data_source = 'app2')
  ) AS best_bmi
FROM raw_patients p JOIN raw_vitalsigns v USING (idcard_hash)
GROUP BY idcard_hash;
\end{lstlisting}

% =============================================================================
\section{App2 Normalizer}
% =============================================================================

App2 เป็นไฟล์เดียวที่แตกต่างจาก Portal/App1 มาก: \textbf{pre-aggregated} (1 row ครอบคลุม
patient + vitalsigns + homevisit + homehealth + lab), ค่าเป็น \textit{Thai labels},
และมี \textit{pre-computed derived} (\texttt{BMI, AGE\_GROUP, LAB\_EGFR}).

\subsection{Split strategy}

1 row ใน \texttt{app2.csv} ถูก split เป็น 5 rows กระจาย 5 ตาราง:

\begin{lstlisting}[style=shellstyle]
app2.csv (1 row)
   |-> raw_patients    (sex, age, age_group_src)
   |-> raw_vitalsigns  (screening flags + bmi_src + *_src labels)
   |-> raw_homevisit   (district + edu_src, occupation_src, ...)
   |-> raw_homehealth  (family history bools + smoke_src, alcohol_src, ...)
   |-> raw_lab_results (binary results + 3 numeric labs + 11 text interpret)
\end{lstlisting}

\subsection{Thai $\rightarrow$ code mapping}

\begin{center}
\footnotesize
\begin{tabularx}{\linewidth}{L{4.5cm} L{4cm} X}
\toprule
\textbf{ค่าดิบ App2} & \textbf{ค่าหลัง normalize} & \textbf{ที่ใช้} \\
\midrule
"ชาย" / "หญิง"                         & 10 / 20                    & \texttt{MALE} $\rightarrow$ \texttt{sex} \\
"ปกติ" / "เสี่ยง" / "เป็น" / "ผิดปกติ"    & 0 / 1 / 1 / 1              & \texttt{DM\_NAME} etc. $\rightarrow$ \texttt{found\_*} \\
"ไม่เป็น" / "เป็น" / "ไม่แน่ใจ"           & False / True / NULL        & \texttt{H\_DM\_NAME} $\rightarrow$ \texttt{parent\_*} \\
"15-34 ปี" / "35-44 ปี" / \ldots        & proxy age 25 / 40 / 52 / 70 & \texttt{AGE\_GROUP} $\rightarrow$ \texttt{age} \\
9999                                   & NULL                        & \texttt{DISTRICT, WRKDISTRICT} \\
Thai label อื่น (e.g., "อ้วนระดับ 1")    & เก็บใน \texttt{*\_src}     & \texttt{BMI\_GROUP, HOMETYPE\_NAME} \\
\bottomrule
\end{tabularx}
\end{center}

\subsection{สิ่งที่ App2 ไม่มี}

\begin{itemize}[leftmargin=*, itemsep=1pt]
  \item \textbf{ไม่มี} \texttt{HEIGHT/WEIGHT} $\rightarrow$ \texttt{bmi} (ETL calc) $=$ NULL,
    ใช้ได้แค่ \texttt{bmi\_src} (App2's).
  \item \textbf{ไม่มี} \texttt{SBP/DBP} $\rightarrow$ \texttt{map\_bp, pulse\_pressure, bp\_group} $=$ NULL.
  \item \textbf{ไม่มี} \texttt{SCN9Q/ST50} individual items $\rightarrow$ \texttt{phq9\_total, st5\_total} $=$ NULL.
  \item \textbf{ไม่มี} \texttt{birth\_year} $\rightarrow$ ใช้ \texttt{AGE\_GROUP} midpoint เป็น proxy.
\end{itemize}

App2 rows จะมี \texttt{NULL} มากกว่า Portal/App1 สำหรับ ETL-derived fields.
การเทียบต้องใช้ \texttt{COALESCE(bmi, bmi\_src)} เพื่อไม่ให้ App2 ตกหล่น.

% =============================================================================
\section{Derived Fields 3 ชั้น}
% =============================================================================

ทุก raw table มี \textbf{3 ชั้น} ของ column:

\subsection{Layer 1 --- Raw}

ค่าคัดจาก CSV ตรง ๆ (หลัง out-of-range $\rightarrow$ NULL).
ตัวอย่าง: \texttt{height\_cm, weight\_kg, sbp, dbp, hemoglobin, cholesterol, fbs}.

\subsection{Layer 2 --- Source-provided (\texttt{*\_src})}

ค่าที่ source system คำนวณไว้ให้ --- เก็บเพื่อ audit.

\begin{center}
\footnotesize
\begin{tabularx}{\linewidth}{l X}
\toprule
\textbf{Column} & \textbf{ที่มา} \\
\midrule
\texttt{bmi\_src}            & App2's pre-computed BMI \\
\texttt{bmi\_group\_src}     & App2's "อ้วนระดับ 1" / "ปกติสมส่วน" / \ldots \\
\texttt{age\_group\_src}     & App2's "15-34 ปี" / \ldots \\
\texttt{cbc\_interp\_src}    & App2's "ปกติ" / "ผิดปกติ" / "ไม่ได้ตรวจ" \\
\texttt{hemoglobin\_src}     & App2's \texttt{LAB\_HEMOGLOBIN} numeric \\
\texttt{edu\_src}            & App2's "ปริญญาตรี" / "มัธยมศึกษา" (ไม่แปลงเป็น code) \\
\bottomrule
\end{tabularx}
\end{center}

\subsection{Layer 3 --- ETL-derived (plain name)}

เราคำนวณเองจาก raw --- \textbf{harmonized ทุก source} (ถ้า raw ครบ).

\begin{center}
\footnotesize
\begin{tabularx}{\linewidth}{l L{5cm} X}
\toprule
\textbf{Column} & \textbf{สูตร} & \textbf{เกณฑ์} \\
\midrule
\texttt{bmi}            & $\dfrac{\text{weight\_kg}}{(\text{height\_cm}/100)^2}$ & $10 \le$ bmi $\le 80$ else NULL \\
\texttt{bmi\_group}     & int 1-5            & WHO Asia-Pacific: ผอม/ปกติ/เกิน/อ้วน I/อ้วน II \\
\texttt{bp\_group}      & int 1-4            & ปกติ / สูงกว่าปกติ / HTN-I / HTN-II \\
\texttt{map\_bp}        & $\text{DBP} + \dfrac{\text{SBP}-\text{DBP}}{3}$ & Mean Arterial Pressure \\
\texttt{pulse\_pressure} & $\text{SBP} - \text{DBP}$                       & --- \\
\texttt{phq9\_total}    & $\sum_{i=1}^{9} Q_i$ (PHQ-9) & NULL ถ้า item ใดว่าง \\
\texttt{phq9\_severity} & 0--4               & ปกติ / เล็ก / กลาง / ค่อนข้าง / รุนแรง \\
\texttt{st5\_total}     & $\sum_{i=1}^{5} Q_i$ (ST-5) & --- \\
\texttt{st5\_severity}  & 0--4               & น้อย / ปานกลาง / สูง / รุนแรง / รุนแรงมาก \\
\texttt{age\_group}     & label              & วัยเรียน / วัยเริ่มทำงาน / \ldots / สูงวัย \\
\texttt{egfr\_stage}    & text               & G1 / G2 / G3a / G3b / G4 / G5 (KDIGO 2012) \\
\texttt{anemia\_class}  & text               & microcytic / normocytic / macrocytic \\
\bottomrule
\end{tabularx}
\end{center}

% =============================================================================
\section{การอัปโหลดผ่านหน้าเว็บ}
% =============================================================================

\subsection{Memory-safe upload}

\begin{center}
\footnotesize
\begin{tabular}{lll}
\toprule
\textbf{ขั้น} & \textbf{เทคนิค} & \textbf{Peak memory} \\
\midrule
Browser $\rightarrow$ Server & multipart streaming            & $< 100$ KB buffer \\
Server receive              & \texttt{SpooledTemporaryFile} (1 MB spool) & $\sim 1$ MB $\times N$ files \\
Copy to tempfile            & \texttt{shutil.copyfileobj} 64 KB chunks    & 64 KB per copy \\
Read into pandas            & \texttt{pd.read\_csv(path)} one-at-a-time   & 1 DataFrame \\
\bottomrule
\end{tabular}
\end{center}

ประมวลผลได้ไฟล์ $> 1$ GB ต่อไฟล์ บนเครื่อง RAM 4 GB.
Max per-file $= 1$ GB (ตั้งใน \texttt{admin.py:MAX\_FILE\_BYTES}).

\subsection{Flow}

\begin{enumerate}[leftmargin=*, itemsep=1pt]
  \item Drag folder $\rightarrow$ JS เก็บ \texttt{file.webkitRelativePath} ใน hidden input.
  \item JS แสดง badge per-file: source (portal/app1/app2/?) + type.
  \item Submit: CSRF check, stream each file to tempfile, build manifest,
    INSERT \texttt{import\_history} (status=\texttt{running}), spawn thread.
  \item Worker: advisory lock $\rightarrow$ per-source delete (ไม่ TRUNCATE)
    $\rightarrow$ loop portal/app1/app2 $\rightarrow$ \texttt{pd.read\_csv} + importer
    $\rightarrow$ commit $\rightarrow$ refresh views $\rightarrow$ flush cache.
  \item UI polls \texttt{/admin/api/import-progress} ทุก 2 วินาที แสดง \texttt{rows\_processed/rows\_total},
    throughput (\texttt{rows/s}), ETA.
\end{enumerate}

% =============================================================================
\section{Schema Changes (Migrations 011--013)}
% =============================================================================

\begin{center}
\footnotesize
\begin{tabularx}{\linewidth}{l X}
\toprule
\textbf{Migration} & \textbf{สิ่งที่ทำ} \\
\midrule
\texttt{011\_multi\_source\_stack.sql}
    & เพิ่ม \texttt{data\_source + import\_batch\_id} ทุก raw table;
      drop \texttt{UNIQUE(idcard\_hash)} $\rightarrow$ \texttt{UNIQUE(idcard\_hash, data\_source)};
      downgrade natural-key UNIQUE $\rightarrow$ plain INDEX;
      เพิ่ม \texttt{*\_src} columns (App2 pre-computed);
      เพิ่ม ETL-derived columns (BMI group, BP group, MAP, pulse pressure,
      PHQ9/ST5 totals + severities, eGFR stage, anemia class);
      สร้าง views \texttt{v\_cross\_system\_duplicates}, \texttt{v\_source\_row\_counts}. \\
\addlinespace
\texttt{012\_matviews\_per\_source.sql}
    & DROP + CREATE matviews ทั้ง 13 views ด้วย \texttt{data\_source} column แรก
      + GROUP BY source. UNIQUE index เป็น \texttt{(data\_source, district\_code, \ldots)}.
      REFRESH ทุก view ท้าย migration. \\
\addlinespace
\texttt{013\_live\_progress.sql}
    & เพิ่ม \texttt{rows\_total + rows\_processed} บน \texttt{import\_history}
      สำหรับ live progress bar (throttled 500 ms per update). \\
\bottomrule
\end{tabularx}
\end{center}

\subsection{ETL ใหม่ --- รองรับ 3 source + progress callback}

\begin{lstlisting}[style=pythonstyle]
# Every importer takes (cur, df, patient_map, data_source, batch_id, progress_cb)
# patient_map key = (idcard_hash, data_source)

def execute_values(cur, sql, rows, page_size=2000, progress_cb=None):
    for chunk in chunks(rows, page_size):
        try:    SAVEPOINT; INSERT chunk; RELEASE
        except: ROLLBACK; retry per-row
        if progress_cb:
            progress_cb(done, total)   # UI update (throttled)
\end{lstlisting}

% =============================================================================
\section{Query Patterns}
% =============================================================================

\subsection{Data quality per source}

\begin{lstlisting}[style=sqlstyle]
-- % null per field per source (for quality comparison)
SELECT data_source,
       COUNT(*) AS total,
       ROUND(100.0 * COUNT(bmi) / COUNT(*), 1) AS pct_bmi,
       ROUND(100.0 * COUNT(sbp) / COUNT(*), 1) AS pct_sbp,
       ROUND(100.0 * COUNT(fbs) / COUNT(*), 1) AS pct_fbs
FROM raw_vitalsigns
GROUP BY data_source;
\end{lstlisting}

\subsection{Source-filtered summary view}

\begin{lstlisting}[style=sqlstyle]
-- หลัง migration 012 ทุก summary_* view มี data_source col
SELECT district_code, SUM(total_screened) AS total_all_sources
FROM summary_district_disease
WHERE data_source IN ('portal', 'app1')   -- Portal + App1 รวม
GROUP BY district_code;
\end{lstlisting}

\subsection{Cross-system duplicate report}

\begin{lstlisting}[style=sqlstyle]
SELECT idcard_hash, n_systems, systems, n_patient_rows
FROM v_cross_system_duplicates
ORDER BY n_systems DESC
LIMIT 100;
-- คนที่ตรวจมากกว่า 1 ระบบ (เช่น portal+app1 = 2, หรือทั้ง 3)
\end{lstlisting}

% =============================================================================
\section{Analytics Dashboards}
% =============================================================================

\subsection{Dashboard / Source tabs}

\texttt{/admin/dashboard?source=portal|app1|app2|all} --- แถบ 4 tab ด้านบน แสดง
per-source counts ของ raw tables และ materialized views.

\subsection{Cross-Stats (3 tabs)}

\texttt{/admin/cross-stats}:

\begin{itemize}[leftmargin=*, itemsep=1pt]
  \item \textbf{Coverage}: Heatmap $50$ เขต $\times$ $3$ source.
    Emerald/sky/amber alpha-scaled ตามจำนวน patients.
  \item \textbf{Distribution}: 9 metrics (\texttt{risk\_dm, risk\_hpt, \ldots, found\_obesity, found\_dyslipidemia})
    แสดง \% per source + $\chi^2$ test แสดง \textit{ต่างกันอย่างมีนัยสำคัญ} ($p < 0.05$) หรือไม่.
    ใช้ 2$\times$N contingency table (positive vs negative $\times$ sources with data).
  \item \textbf{Data Quality}: $230$ fields $\times$ $7$ tables field-completeness $\%$.
    Blank cell (\texttt{—}) สำหรับ source ที่ไม่มีข้อมูลใน column นั้น.
\end{itemize}

\subsection{Agreement (Phase 3) --- Same-person discrepancy}

\texttt{/admin/agreement?pair=portal-app1|portal-app2|app1-app2}.

Match ผู้ป่วยข้าม source ด้วย \texttt{idcard\_hash} (DISTINCT ON + ORDER BY visit\_date DESC
$\rightarrow$ latest visit per source). จากนั้น:

\begin{itemize}[leftmargin=*, itemsep=1pt]
  \item \textbf{Bland-Altman} (continuous: BMI, SBP, DBP, Hb, Cholesterol, eGFR, \ldots):
    bias, SD, $95\%$ LoA $=$ bias $\pm 1.96 \cdot$ SD, Pearson $r$, ICC.
    Matplotlib scatter plot embedded as base64 PNG.
  \item \textbf{Cohen's $\kappa$} (categorical: DM/HPT/CVD/obesity found, risk flags, sex):
    confusion matrix $+$ Landis-Koch interpretation
    (slight / fair / moderate / substantial / almost perfect).
\end{itemize}

$N \ge 3$ คู่ ถึงจะคำนวณ, $N \ge 10$ ถึงจะวาด plot. ต่ำกว่านั้นแสดงข้อความเตือนแทน.

% =============================================================================
\section{Troubleshooting}
% =============================================================================

\begin{center}
\footnotesize
\begin{tabularx}{\linewidth}{L{5cm} X}
\toprule
\textbf{Issue} & \textbf{Fix} \\
\midrule
\texttt{IDCARD\_HASH\_SECRET} ไม่ set
    & \texttt{export IDCARD\_HASH\_SECRET=\$(python3 -c 'import secrets; print(secrets.token\_urlsafe(32))')} \\
\addlinespace
"Another import is currently running"
    & PostgreSQL advisory lock --- รอ import เก่าเสร็จ หรือ restart API เพื่อ drop dead connection \\
\addlinespace
Import status "success" แต่ rows น้อย
    & ดู \texttt{rows\_skipped} + ตรวจ log --- มักเพราะ PID decode ไม่ได้ \\
\addlinespace
View refresh failed
    & \texttt{import\_history.view\_refresh\_error};
      re-run ผ่าน \texttt{POST /admin/refresh} หรือ \texttt{make refresh-views} \\
\addlinespace
Import thread orphaned (API restart)
    & \texttt{UPDATE import\_history SET status='error' WHERE status='running';} \\
\addlinespace
Memory error ระหว่าง import
    & DataFrame too large --- ลด file size หรือเพิ่ม Docker memory limit \\
\bottomrule
\end{tabularx}
\end{center}

% =============================================================================
\section{ไฟล์ที่เกี่ยวข้อง}
% =============================================================================

\begin{center}
\footnotesize
\begin{tabularx}{\linewidth}{l X}
\toprule
\textbf{Path} & \textbf{หน้าที่} \\
\midrule
\texttt{db/migrations/011\_multi\_source\_stack.sql} & เปลี่ยน schema สำหรับ 3-source stack mode \\
\texttt{db/migrations/012\_matviews\_per\_source.sql} & Rebuild 13 matviews + data\_source GROUP BY \\
\texttt{db/migrations/013\_live\_progress.sql} & เพิ่ม rows\_total/rows\_processed \\
\texttt{etl/import\_csv.py} & Main ETL orchestrator + per-table importers + progress cb \\
\texttt{etl/derived.py} & Shared derived helpers (BMI, MAP, PHQ9, eGFR, anemia) \\
\texttt{etl/app2\_normalizer.py} & App2 Thai$\rightarrow$code mapping + row splitter \\
\texttt{api/admin.py} & Bundle upload + background worker + cross-stats + agreement \\
\texttt{api/services/agreement\_service.py} & Bland-Altman + Cohen's $\kappa$ + plot renderer \\
\texttt{api/templates/admin/upload\_bundle.html} & Upload UI (folder picker) \\
\texttt{api/templates/admin/cross\_stats.html} & Coverage heatmap + distribution + data quality \\
\texttt{api/templates/admin/agreement.html} & Agreement report UI \\
\texttt{PREPROCESSING.md} & เอกสาร Markdown (ต้นฉบับของเอกสารนี้) \\
\texttt{DATA-DICTIONARY.md} & Column-level reference \\
\texttt{MEDICAL-DICTIONARY.md} & Clinical range reference \\
\bottomrule
\end{tabularx}
\end{center}

\vspace{0.5cm}
\begin{center}
\small\color{gray}\itshape
เอกสารสร้างเมื่อ 2026-04-22 · ครอบคลุมถึง migration 013 + Cross-Stats/Agreement dashboards
\end{center}

\end{document}
